\section{Đánh giá một lời giải thích}
\label{sec:ch01-danh-gia-giai-thich}

Mục~\ref{sec:ch01-khi-du-doan-chua-du} đã cho thấy mỗi mục đích đòi một việc
kiểm tra riêng. Nếu mục đích quyết định thước đo thì không thể đánh giá mọi
lời giải thích bằng một con số hoặc một đại lượng thay thế (proxy) mang tính
cấu trúc. Doshi-Velez và Kim đề xuất ba cách đặt bài toán đánh giá, khác nhau
ở người tham gia và nhiệm vụ được dùng~\autocite{p14rigorous}.
Hình~\ref{fig:ch01-ba-cach-danh-gia} đặt ba cách ấy trên cùng một trục.

\begin{figure}[htbp]
  \centering
  \begin{tikzpicture}[
  box/.style={draw=black!65, rounded corners=2pt, align=center,
    minimum width=4.4cm, minimum height=1.1cm, fill=black!5},
  data/.style={draw=black!65, rounded corners=2pt, align=center,
    minimum width=4.4cm, minimum height=1.1cm, fill=black!15},
  arrow/.style={-{Stealth[length=2mm]}, draw=black!65, thick}
]
  \node[data] (application) {Đánh giá trong ứng dụng};
  \node[box, below=18mm of application] (human) {Đánh giá với người dùng};
  \node[box, below=18mm of human] (functional) {Đánh giá theo chức năng};

  \draw[arrow] (human) -- (application);
  \draw[arrow] (functional) -- (human);

  \node[align=left, anchor=west] at ($(human.north)!0.5!(application.south) + (8mm,0)$)
    {vay giả định rằng nhiệm vụ rút gọn\\giữ được phần cốt lõi của nhiệm vụ thật};
  \node[align=left, anchor=west] at ($(functional.north)!0.5!(human.south) + (8mm,0)$)
    {vay giả định rằng đại lượng thay thế\\đã được kiểm chứng với người dùng};
\end{tikzpicture}

  \caption{Ba cách đánh giá khác nhau ở bằng chứng mà chúng đòi hỏi. Đánh giá
  trong ứng dụng dùng người dùng và nhiệm vụ thực nên gần mục tiêu triển khai
  nhất. Đánh giá với người dùng rút gọn nhiệm vụ để cô lập một giả thuyết.
  Đánh giá theo chức năng dùng proxy, rẻ hơn nhưng phải biện minh vì sao proxy
  có liên hệ với điều người dùng cần. Sự tiện lợi của cách rẻ hơn không chuyển
  thành bằng chứng cho cách gần mục tiêu hơn.}
  \label{fig:ch01-ba-cach-danh-gia}
\end{figure}

\paragraph{Đánh giá trong ứng dụng.} Người dùng thật thực hiện nhiệm vụ thật;
ví dụ, một bác sĩ dùng lời giải thích để phát hiện khi mô hình dự báo dựa vào
một tín hiệu không phù hợp. Đây là cách đánh giá gần với mục đích triển khai
nhất, nhưng cũng tốn kém và khó kiểm soát các yếu tố nhiễu.

\paragraph{Đánh giá với người dùng.} Người tham gia vẫn là con người, nhưng
nhiệm vụ được giản lược; chẳng hạn, ta có thể yêu cầu họ chọn lời giải thích
nào giúp họ dự đoán hành vi mô hình tốt hơn. Thiết kế này không tái tạo toàn
bộ công việc thực tế, nhưng cho phép kiểm tra một thuộc tính xác định trong
cách trình bày lời giải thích mà không phải dựng cả ứng dụng.

\paragraph{Đánh giá theo chức năng.} Cách thứ ba không cần người tham gia:
nghiên cứu đo một proxy như độ thưa của mô hình, độ ngắn của quy tắc, hay khả
năng dự đoán một nhãn thay thế. Cách làm này phù hợp khi proxy đã có cơ sở,
nhưng proxy không phải bản thân mục tiêu. Nghiên cứu so sánh các mô hình dễ
diễn giải gần đây cũng phải chọn proxy cấu trúc để có thể so sánh trên quy mô
lớn, và thừa nhận giới hạn của lựa chọn ấy~\autocite{p15microscope}.

Như vậy, ba cách trên không phải là ba nhãn chất lượng cho cùng một đại
lượng, mà là ba câu trả lời cho ba câu hỏi thực nghiệm khác nhau. Một kết quả
tốt ở cách đánh giá theo chức năng có thể là lý do để làm nghiên cứu tiếp,
không phải giấy chứng nhận rằng người dùng sẽ hiểu đúng hoặc rằng lời giải
thích đã trung thực.
